% GENERATED FILE. Edit canonical lessons, then run npm run generate:books.
% canonical-source-hash: fnv1a64:18a66e3ef666872f
% canonical-lessons: ZH-C46-xingqi, ZH-W46-xing, ZH-W46-qi-period, ZH-C46-xingqiyi, ZH-C46-shihou, ZH-C46-zhoumo, ZH-C46-xingqitian

\chapter{The Week}
\label{ch:zh-the-week}
\hlchaptermodality{46}

\begin{chapteropening}
\textbf{By the end of this chapter:} \emph{I can name the week, Monday, Sunday and the weekend.}

Complete the last lesson of chapter 46: I can name the week, Monday, Sunday and the weekend.
\end{chapteropening}

\section[\texorpdfstring{\zh{星期} (xīngqī)}{星期 (xīngqī)}]{\zh{星期} (xīngqī) — week}

\label{lesson:ZH-C46-xingqi}



\begin{quote}
Before the new one: say the Mandarin for now, then the Mandarin for tomorrow.
\end{quote}

\subsection*{You'll want to know: \zh{星期}}
\textbf{\zh{星期}} — \emph{xīngqī} — \textquotedblleft{}week\textquotedblright{}.

\begin{quote}
\textbf{\zh{一个星期有七天}} \emph{yí ge xīngqī yǒu qī tiān} — There are seven days in a week.
\end{quote}

The new character, in pieces: \zh{星} — \zh{日} above \zh{生}; \zh{期} — an eight-stroke piece on the left and a moon piece on the right. The next lesson puts it in your hand.

\begin{cousinweb}
\textbf{\zh{星期}} (\emph{xīngqī}) is built from \zh{星} 'star' and \zh{期} 'appointed time', both inherited from Middle Chinese, \emph{seng} and \emph{gi} (Guangyun, in Baxter's transcription). How the pair came to mean a week is not settled.
\end{cousinweb}

\begin{grammarlens}[title={What you've built}]
One more word for times and partings.
\end{grammarlens}

\subsection*{Your turn}
\begin{itemize}
\raggedright
  \item \emph{Say it:} \emph{xīngqī}
  \item \emph{Say it:} \emph{xīngqī}, once more
  \item \emph{Say it:} \emph{yí ge xīngqī yǒu qī tiān} — There are seven days in a week.
  \item \emph{Recall:} say the Mandarin for now, then the Mandarin for tomorrow, then say \emph{xīngqī} again
\end{itemize}

\begin{tcolorbox}[breakable,colback=teal!4,colframe=teal!35!black,title={Before you move on}]
What does \zh{星期} mean? (\textquotedblleft{}Week\textquotedblright{}.) Say it once more.
\end{tcolorbox}
\section[\texorpdfstring{\zh{星} (xīng)}{星 (xīng)}]{\zh{星} — \zh{日} above \zh{生}}

\label{lesson:ZH-W46-xing}



\begin{quote}
The word first, then the character inside it.

\begin{itemize}
\raggedright
  \item \emph{Recall:} say \emph{xīngqī} — \textbf{R1}, one lesson back, before it decays
\end{itemize}

\begin{quote}
\zh{日} \zh{生}
\end{quote}

\emph{Read it:} each one aloud, with nothing said first — \emph{rì}, \emph{shēng}
\end{quote}

\subsection*{Script — \zh{日} above \zh{生}}
\hlblockfigure{\detokenize{figures/ZH-W46-xing-filmstrip.pdf}}{How \zh{星} is written, stroke by stroke}

\begin{quote}
\zh{星}
\end{quote}

\textbf{Two pieces you know, top to bottom: \zh{日}, then \zh{生}.}

\textbf{\zh{日}, four strokes}, short and wide on top.

\textbf{\zh{生}, five strokes}, as in \zh{学生}: a left-falling stroke, a horizontal, a shorter horizontal, the vertical through them, and the long bar at the foot.

4 + 5 = nine.

\subsection*{Your turn}
\begin{itemize}
\raggedright
  \item \emph{Write it:} \zh{星}, all nine, with the model in front of you
  \item \emph{Cover:} write it again from memory, and count the pieces, not the strokes
\end{itemize}

\begin{tcolorbox}[breakable,colback=teal!4,colframe=teal!35!black,title={Before you move on}]
Which two characters you know make \zh{星}? (\textbf{\zh{日} and \zh{生}}.) Which is on top? (\textbf{\zh{日}}.) Which stroke of \zh{生} is last? (\textbf{The long bar at the foot}.)

Source: \href{https://raw.githubusercontent.com/chanind/hanzi-writer-data/68d10a4b21150cae5e1ebbd223eed289cf32d90c/data/\%E6\%98\%9F.json}{Hanzi Writer Data \zh{星}, pinned stroke paths and medians}.
\end{tcolorbox}
\section[\texorpdfstring{\zh{期} (qī)}{期 (qī)}]{\zh{期} — eight strokes, then the moon}

\label{lesson:ZH-W46-qi-period}



\begin{quote}
The word first, then the character inside it.

\begin{itemize}
\raggedright
  \item \emph{Recall:} say \emph{xīngqī} — \textbf{R1}, one lesson back, before it decays
\end{itemize}

\begin{quote}
\zh{有} \zh{日}
\end{quote}

\emph{Read it:} each one aloud, with nothing said first — \emph{yǒu}, \emph{rì}
\end{quote}

\subsection*{Script — eight strokes, then the moon}
\hlblockfigure{\detokenize{figures/ZH-W46-qi-period-filmstrip.pdf}}{How \zh{期} is written, stroke by stroke}

\begin{quote}
\zh{期}
\end{quote}

\textbf{Two pieces, left before right.}

\textbf{The left piece, eight strokes.} A horizontal at the top. Two verticals down through it. Two short horizontals between them. A long horizontal across the foot. Then a left-falling stroke and a dot under it.

\textbf{The moon, four strokes.} A left-falling stroke down the left. A horizontal that turns down the right and hooks. Two short horizontals inside. It is the lower part of \zh{有}.

8 + 4 = twelve.

\subsection*{Your turn}
\begin{itemize}
\raggedright
  \item \emph{Write it:} \zh{期}, all twelve, with the model in front of you
  \item \emph{Cover:} write it again from memory, and count the pieces, not the strokes
\end{itemize}

\begin{tcolorbox}[breakable,colback=teal!4,colframe=teal!35!black,title={Before you move on}]
How many pieces in \zh{期}? (\textbf{Two} — eight strokes, then the moon.) What is under the long bar of the left piece? (\textbf{A left-falling stroke and a dot}.) How does the moon begin? (\textbf{With a left-falling stroke}.)

Source: \href{https://raw.githubusercontent.com/chanind/hanzi-writer-data/68d10a4b21150cae5e1ebbd223eed289cf32d90c/data/\%E6\%9C\%9F.json}{Hanzi Writer Data \zh{期}, pinned stroke paths and medians}.
\end{tcolorbox}
\section[\texorpdfstring{\zh{星期一} (xīngqīyī)}{星期一 (xīngqīyī)}]{\zh{星期一} (xīngqīyī) — Monday}

\label{lesson:ZH-C46-xingqiyi}



\begin{quote}
Before the new one: say the Mandarin for tomorrow, then the Mandarin for week.
\end{quote}

\subsection*{You'll want to know: \zh{星期一}}
\textbf{\zh{星期一}} — \emph{xīngqīyī} — \textquotedblleft{}Monday\textquotedblright{}.

\begin{quote}
\textbf{\zh{今天星期一}} \emph{jīntiān xīngqīyī} — Today's Monday.
\end{quote}

Tone: \zh{一} keeps its first tone here, \emph{yī}. It counts a day in order, and the change you know from \emph{yìbǎi} does not apply.

\begin{cousinweb}
\textbf{\zh{星期一}} (\emph{xīngqīyī}) is built from \zh{星期} 'week' and \zh{一} 'one': day one of the week.
\end{cousinweb}

\begin{grammarlens}[title={What you've built}]
One more word for times and partings.
\end{grammarlens}

\subsection*{Your turn}
\begin{itemize}
\raggedright
  \item \emph{Say it:} \emph{xīngqīyī}
  \item \emph{Say it:} \emph{xīngqīyī}, once more
  \item \emph{Say it:} \emph{jīntiān xīngqīyī} — Today's Monday.
  \item \emph{Recall:} say the Mandarin for tomorrow, then the Mandarin for week, then say \emph{xīngqīyī} again
\end{itemize}

\begin{tcolorbox}[breakable,colback=teal!4,colframe=teal!35!black,title={Before you move on}]
What does \zh{星期一} mean? (\textquotedblleft{}Monday\textquotedblright{}.) Say it once more.
\end{tcolorbox}
\section[shíhou]{shíhou — time; moment}

\label{lesson:ZH-C46-shihou}



\begin{quote}
Before the new one: say the Mandarin for week, then the Mandarin for Monday.
\end{quote}

\subsection*{You'll want to know: shíhou}
\emph{shíhou} — \textquotedblleft{}time; moment\textquotedblright{}.

\begin{quote}
\emph{nǐ shénme shíhou qù} — When are you going?
\end{quote}

Its characters come later; for now it is a word for the ear.

\begin{cousinweb}
\emph{shíhou} is built from shí 'time' and hou 'to wait for', both inherited from Middle Chinese, \emph{dzyi} and \emph{huwH} (Guangyun, in Baxter's transcription).
\end{cousinweb}

\begin{grammarlens}[title={What you've built}]
One more word for times and partings.
\end{grammarlens}

\subsection*{Your turn}
\begin{itemize}
\raggedright
  \item \emph{Say it:} \emph{shíhou}
  \item \emph{Say it:} \emph{shíhou}, once more
  \item \emph{Say it:} \emph{nǐ shénme shíhou qù} — When are you going?
  \item \emph{Recall:} say the Mandarin for week, then the Mandarin for Monday, then say \emph{shíhou} again
\end{itemize}

\begin{tcolorbox}[breakable,colback=teal!4,colframe=teal!35!black,title={Before you move on}]
What does shíhou mean? (\textquotedblleft{}Time; moment\textquotedblright{}.) Say it once more.
\end{tcolorbox}
\section[zhōumò]{zhōumò — weekend}

\label{lesson:ZH-C46-zhoumo}



\begin{quote}
Before the new one: say the Mandarin for Monday, then the Mandarin for moment.
\end{quote}

\subsection*{You'll want to know: zhōumò}
\emph{zhōumò} — \textquotedblleft{}weekend\textquotedblright{}.

\begin{quote}
\emph{zhōumò wǒ xǐhuan kàn shū} — At the weekend I like to read.
\end{quote}

Its characters come later; for now it is a word for the ear.

\begin{cousinweb}
\emph{zhōumò} is built from zhōu 'week' and mò 'end'. The second is inherited from Middle Chinese \emph{mat} (Guangyun, in Baxter's transcription), glossed 'the tip of a tree' and 'end'. The first character is recorded there as \emph{tsyuw}, 'a full circuit'; the sense week is not given, so how it came to name the week is not settled.
\end{cousinweb}

\begin{grammarlens}[title={What you've built}]
One more word for times and partings.
\end{grammarlens}

\subsection*{Your turn}
\begin{itemize}
\raggedright
  \item \emph{Say it:} \emph{zhōumò}
  \item \emph{Say it:} \emph{zhōumò}, once more
  \item \emph{Say it:} \emph{zhōumò wǒ xǐhuan kàn shū} — At the weekend I like to read.
  \item \emph{Recall:} say the Mandarin for Monday, then the Mandarin for moment, then say \emph{zhōumò} again
\end{itemize}

\begin{tcolorbox}[breakable,colback=teal!4,colframe=teal!35!black,title={Before you move on}]
What does zhōumò mean? (\textquotedblleft{}Weekend\textquotedblright{}.) Say it once more.
\end{tcolorbox}
\section[\texorpdfstring{\zh{星期天} (xīngqītiān)}{星期天 (xīngqītiān)}]{\zh{星期天} (xīngqītiān) — Sunday}

\label{lesson:ZH-C46-xingqitian}



\begin{quote}
Before the new one: say the Mandarin for moment, then the Mandarin for weekend.
\end{quote}

\subsection*{You'll want to know: \zh{星期天}}
\textbf{\zh{星期天}} — \emph{xīngqītiān} — \textquotedblleft{}Sunday\textquotedblright{}.

\begin{quote}
\textbf{\zh{我喜欢星期天}} \emph{wǒ xǐhuan xīngqītiān} — Sunday is the day I like best.
\end{quote}

\begin{cousinweb}
\textbf{\zh{星期天}} (\emph{xīngqītiān}) is built from \zh{星期} 'week' and \zh{天} 'day'.
\end{cousinweb}

\begin{grammarlens}[title={What you've built}]
One more word for times and partings.
\end{grammarlens}

\subsection*{Your turn}
\begin{itemize}
\raggedright
  \item \emph{Say it:} \emph{xīngqītiān}
  \item \emph{Say it:} \emph{xīngqītiān}, once more
  \item \emph{Say it:} \emph{wǒ xǐhuan xīngqītiān} — Sunday is the day I like best.
  \item \emph{Recall:} say the Mandarin for moment, then the Mandarin for weekend, then say \emph{xīngqītiān} again
\end{itemize}

\begin{tcolorbox}[breakable,colback=teal!4,colframe=teal!35!black,title={Before you move on}]
What does \zh{星期天} mean? (\textquotedblleft{}Sunday\textquotedblright{}.) Say it once more.
\end{tcolorbox}
